\documentclass[10pt, aspectratio=169]{beamer}

% --- PAKIETY I KONFIGURACJA ---
\usepackage[utf8]{inputenc}
\usepackage[polish]{babel}
\usepackage{polski}
\usepackage{helvet}
\usepackage{tikz}
\usepackage{booktabs}
\usepackage[most]{tcolorbox}

% --- PALETA BARW ---
\definecolor{bgCanvas}{RGB}{250, 252, 255}
\definecolor{textPrimary}{RGB}{30, 41, 59}
\definecolor{textMuted}{RGB}{100, 116, 139}
\definecolor{primaryNavy}{RGB}{15, 23, 42}
\definecolor{accentCoral}{RGB}{225, 29, 72}
\definecolor{accentTeal}{RGB}{13, 148, 136}
\definecolor{cardBg}{RGB}{255, 255, 255}
\definecolor{cardBorder}{RGB}{226, 232, 240}

% --- KONFIGURACJA STYLU BEAMER ---
\setbeamercolor{background canvas}{bg=bgCanvas}
\setbeamercolor{normal text}{fg=textPrimary}
\setbeamercolor{title}{fg=primaryNavy}
\setbeamercolor{author}{fg=textMuted}
\setbeamercolor{date}{fg=textMuted}
\setbeamercolor{frametitle}{fg=primaryNavy, bg=}

\setbeamertemplate{navigation symbols}{}

\setbeamertemplate{frametitle}{
    \vspace*{0.3cm}
    \begin{beamercolorbox}[wd=\paperwidth, ht=1.0cm, dp=0.2cm, leftskip=0.8cm]{frametitle}
        {\usebeamerfont{frametitle}\Huge\textbf{\insertframetitle}}
    \end{beamercolorbox}
}

\setbeamertemplate{footline}{
    \begin{beamercolorbox}[wd=\paperwidth, ht=0.6cm, dp=0.3cm, rightskip=0.8cm]{footline}
        \hfill \small\textcolor{textMuted}{\insertframenumber\,/\,\inserttotalframenumber}
    \end{beamercolorbox}
}

% --- STYLIZACJA BLOKÓW (TCOLORBOX) ---
\tcbset{
    moderncard/.style={
        enhanced,
        colback=cardBg,
        colframe=cardBorder,
        boxrule=1pt,
        arc=4pt,
        boxsep=5pt,
        left=8pt, right=8pt, top=10pt, bottom=8pt,
        fonttitle=\bfseries\sffamily,
        coltitle=primaryNavy,
        colbacktitle=cardBg,
        attach boxed title to top left={xshift=10pt, yshift=-0.25mm},
        boxed title style={
            boxrule=0pt,
            colframe=cardBg,
            arc=0pt,
            outer arc=0pt,
            top=2pt, bottom=2pt, left=4pt, right=4pt
        }
    }
}

% --- DANE PREZENTACJI ---
\title{Publiczne ocenianie, porównywanie i zawstydzanie uczniów oraz ich konsekwencje}
\subtitle{Sztuka budowania bezpiecznego środowiska uczenia się}
\author{Prezentacja studencka}
\institute{Kierunek: Matematyka Nauczycielska \quad|\quad Przedmiot: Komunikacja Interpersonalna}
\date{\today}

\begin{document}

% --- SLAJD 1: TYTUŁOWY ---
\begin{frame}[plain]
    \begin{tikzpicture}[remember picture, overlay]
        \fill[primaryNavy] (current page.north west) rectangle ([xshift=2.2cm]current page.south west);
        \fill[accentCoral] ([xshift=2.0cm]current page.north west) rectangle ([xshift=2.2cm]current page.south west);
    \end{tikzpicture}
    
    \vfill
    \begin{beamercolorbox}[wd=\paperwidth, leftskip=3.2cm, rightskip=1.0cm]{title}
        \usebeamerfont{title}
        \sffamily 
        
        {\Large \sffamily \bfseries \textcolor{textMuted}{}}\\[0.4cm]
        
        {\fontsize{26pt}{32pt}\selectfont \sffamily \bfseries \textcolor{primaryNavy}{Publiczne ocenianie,}}\\[0.15cm]
        {\fontsize{26pt}{32pt}\selectfont \sffamily \bfseries \textcolor{primaryNavy}{porównywanie i zawstydzanie}}\\[0.15cm]
        {\fontsize{26pt}{32pt}\selectfont \sffamily \bfseries \textcolor{primaryNavy}{uczniów przez nauczyciela}}\\[0.35cm]
        
        {\fontsize{26pt}{32pt}\selectfont \sffamily \bfseries \textcolor{accentCoral}{i jego konsekwencje}}
    \end{beamercolorbox}
    \vfill
\end{frame}

% --- SLAJD 2: SZKOLNA CODZIENNOŚĆ ---
{
\usebackgroundtemplate{
    \tikz[remember picture, overlay]
    \node[opacity=0.50, inner sep=0pt] at (current page.center)
    {\includegraphics[width=\paperwidth, height=\paperheight]{ooolowek.png}};
}
\begin{frame}{}
    \begin{columns}[T]
        \begin{column}{0.62\textwidth}
            \vspace{0.1cm}
            {\large\bfseries\textcolor{primaryNavy}{Z czym wciąż spotykamy się w klasach?}}
            
            \vspace{0.4cm}
            \begin{itemize}
                \setlength{\itemsep}{12pt} 
                \item \textbf{głośne odczytywanie wyników} sprawdzianów na forum klasy
                \item \textbf{porównywanie uczniów ze sobą} i wyróżnianie jednych kosztem pozostałych
                \item \textbf{krytykowanie uczniów przy tablicy} z powodu braków w wiedzy lub trudności z zadaniem
            \end{itemize}
        \end{column}
\pause
        \begin{column}{0.35\textwidth}
            \begin{tcolorbox}[
                enhanced,
                colback=primaryNavy!5,
                colframe=primaryNavy!20,
                arc=4mm,
                boxrule=0.5pt,
                top=4mm, bottom=4mm, left=4mm, right=4mm
            ]
                \centering
                \small\itshape\textcolor{primaryNavy}{
                    Czy publiczne wytykanie błędów i presja grupy rzeczywiście motywują do nauki...?
                }
            \end{tcolorbox}
            
            \vspace{2.5cm} 
        \end{column}
    \end{columns}
\end{frame}
}

% --- SLAJD 3: MECHANIZMY PRESJI SPOŁECZNEJ ---
{
\usebackgroundtemplate{
    \tikz[remember picture, overlay]
    \node[opacity=0.50, inner sep=0pt] at (current page.center)
    {\includegraphics[width=\paperwidth, height=\paperheight]{kartka_x.jpg}};
}
\begin{frame}{Presja społeczna w klasie}
    \begin{columns}[T]
        \begin{column}{0.20\textwidth}
        \end{column}
        \pause
        \begin{column}{0.65\textwidth}
            \vspace{0.64cm}
            {\large\bfseries\textcolor{primaryNavy}{Publiczna Ocena}}
            \vspace{0.07cm}
            \begin{itemize}
                \footnotesize
                \item Czytanie wyników na głos
                \item Sprowadzanie jednostki wyłącznie \\do cyfry
            \end{itemize}
            \pause
            \vspace{0.3cm}
            
            {\large\bfseries\textcolor{primaryNavy}{Porównywanie}}
            \vspace{0.07cm}
            \begin{itemize}
                \footnotesize
                \item Nadawanie etykiet (prymus vs słaby)
                \item Zaburzenie relacji rówieśniczych \\(rywalizacja)
            \end{itemize}
            \pause
            \vspace{0.3cm}
            
            {\large\bfseries\textcolor{primaryNavy}{Zawstydzanie}}
            \vspace{0.07cm}
            \begin{itemize}
                \footnotesize
                \item Sarkazm i wytykanie błędów
                \item Komunikat: „Jesteś niewystarczający”
            \end{itemize}
            
            \vspace{0.4cm}
            
        \end{column}
        
        \begin{column}{0.225\textwidth}
        \end{column}
    \end{columns}
\end{frame}
}

% --- SLAJD 4: TOKSYCZNOŚĆ PORÓWNYWANIA ---
{
\usebackgroundtemplate{
    \tikz[remember picture, overlay]
    \node[opacity=0.50, inner sep=0pt] at (current page.center)
    {\includegraphics[width=\paperwidth, height=\paperheight]{prymus.jpg}};
}
\begin{frame}{}
    \vspace{0.4cm}
    \begin{columns}[T]
        \begin{column}{0.39\textwidth}
            
        \end{column}
        
        \begin{column}{0.04\textwidth}
        \end{column}
        
        \begin{column}{0.57\textwidth}
            {\Large\sffamily\bfseries\textcolor{primaryNavy}{Negatywne skutki porównywania uczniów ze sobą:}}\\[0.5cm]
            \small
            \begin{itemize}
                \setlength{\itemsep}{10pt}
                \pause\item Publiczne stawianie za wzór jest dla ucznia kłopotliwe – rodzi zawiść rówieśników i grozi wykluczeniem z grupy.
                \pause\item Przy spadku samooceny uczeń zużywa energię \\na obronę własnej godności, zamiast na efektywną naukę.
                \pause\item Rywalizacja uczy postawy \textit{„jeśli inni wypadną źle, ja wydam się lepszy”}. Niszczy to szczerość i chęć do współpracy.
            \end{itemize}
            \vspace{0.4cm}
            \scriptsize\textcolor{textMuted}{Źródło: Th. Gordon, B. Niemierko}
        \end{column}
    \end{columns}
\end{frame}
}

% --- SLAJD 5: NEUROBIOLOGIA ---
{
\usebackgroundtemplate{
    \tikz[remember picture, overlay]
    \node[anchor=west, opacity=0.45, inner sep=0pt] at (current page.west) 
    {\includegraphics[height=\paperheight, keepaspectratio]{mozg2.jpg}};
}
\begin{frame}{Zawstydzanie powoduje stres}
    \pause
    
    \begin{tikzpicture}[remember picture, overlay]
        \node[anchor=south west, text width=0.48\paperwidth, align=left, inner sep=0pt] 
        at ([xshift=1.2cm, yshift=1.2cm]current page.south west) {
            \footnotesize
            
        };
    \end{tikzpicture}

    \vspace{-0.7cm}
    \begin{columns}[T]
        \begin{column}{0.42\textwidth}
        \end{column}
        
        \begin{column}{0.58\textwidth}
            \begin{tcolorbox}[
                enhanced,
                colback=white,
                opacityback=0.95, 
                colframe=cardBorder,
                boxrule=1pt,
                arc=4pt,
                boxsep=5pt,
                left=10pt, right=10pt, top=12pt, bottom=12pt
            ]
                {\large \bfseries \sffamily \textcolor{primaryNavy}{Jak stres blokuje proces poznawczy?}}\\[0.04cm]
                
                \small
                \begin{itemize}
                    \setlength{\itemsep}{10pt}
                    \item Zawstydzenie działa na ciało migdałowate jak fizyczne zagrożenie.
                    \item Kortyzol i adrenalina wymuszają reakcję obronną, tryb przetrwania.
                    \item Kora przedczołowa (logika, analiza i pamięć) zostaje czasowo zablokowana.
                    
                \end{itemize}
            \end{tcolorbox}
        \end{column}
    \end{columns}
\end{frame}
}

% --- SLAJD 6: BLOKADY KOMUNIKACYJNE ---
{
\usebackgroundtemplate{
    \tikz[remember picture, overlay]
    \node[opacity=0.50, inner sep=0pt] at (current page.center)
    {\includegraphics[width=\paperwidth, height=\paperheight]{mur.jpg}};
}
\vspace{0.2cm}
\begin{frame}{\normalsize\textcolor{textPrimary}{
                Typowe automatyczne reakcje, które uniemożliwiają efektywny dialog i \textbf{budują mur} między nauczycielem a uczniem:
            }}
    \begin{columns}[T]
        \begin{column}{0.20\textwidth}
        \end{column}
        
        \begin{column}{0.55\textwidth}
            \pause
            \begin{flushleft}
                {\small\bfseries\textcolor{primaryNavy}{Moralizowanie}}\\
                \textit{\footnotesize „Powinieneś bardziej się postarać, z takim podejściem nic nie osiągniesz.”}
            \end{flushleft}
            \vspace{-0.1cm}
            \pause
            \begin{flushleft}
                \hspace{0.77cm}{\small\bfseries\textcolor{primaryNavy}{Osądzanie}}\\
                \hspace{0.77cm}\textit{\footnotesize „Zawsze popełniasz te same błędy.”}
            \end{flushleft}
            \vspace{-0.1cm}
            \pause
            \begin{flushleft}
                {\small\bfseries\textcolor{primaryNavy}{Ośmieszanie}}\\
                \textit{\footnotesize „To było w szkole podstawowej, a ty wciąż tego \\nie rozumiesz?”}
            \end{flushleft}
            \vspace{-0.1cm}
            \pause
            \begin{flushleft}
                \hspace{0.25cm}{\small\bfseries\textcolor{primaryNavy}{Diagnozowanie}}\\
                \hspace{0.25cm}\textit{\footnotesize „Robisz to specjalnie, żeby przedłużyć lekcję.”}
            \end{flushleft}
            
        \end{column}
        
        \begin{column}{0.43\textwidth}
        \end{column}
    \end{columns}

    \begin{tikzpicture}[remember picture, overlay]
        \pause\node[rotate=15, text=black] at ([xshift=4.5cm, yshift=0.6cm]current page.center) {\scriptsize\bfseries POSTAWA OBRONNA};
        \node[rotate=15, text=black] at ([xshift=4.5cm, yshift=-0.6cm]current page.center) {\scriptsize\bfseries GNIEW};
        \node[rotate=15, text=black] at ([xshift=5.0cm, yshift=-1.5cm]current page.center) {\scriptsize\bfseries WYCOFANIE Z AKTYWNOŚCI};
    \end{tikzpicture}
\end{frame}
}

% --- SLAJD 7: PARADOKS NAUCZYCIELA ---
\begin{frame}{Paradoks nauczyciela}
    \vspace{0.4cm}
    \begin{columns}[T]
        \begin{column}{0.05\textwidth}
        \end{column}
        
        \begin{column}{0.24\textwidth}
            {\fontsize{40}{40}\selectfont\sffamily\bfseries\textcolor{primaryNavy!10}{1}}\\[-0.7cm]
            {\large\bfseries\textcolor{primaryNavy}{Intencja}}\\[0.3cm]
            \small\textcolor{textPrimary}{Chcę przekazać wiedzę, zmotywować klasę i utrzymać dyscyplinę.}
        \end{column}
        
        \begin{column}{0.09\textwidth}
            \vspace{1.2cm}
            \centering\textcolor{textMuted}{\Large $\rightarrow$}
        \end{column}
        
        \begin{column}{0.24\textwidth}
            {\fontsize{40}{40}\selectfont\sffamily\bfseries\textcolor{accentCoral!10}{2}}\\[-0.7cm]
            {\large\bfseries\textcolor{accentCoral}{Metoda}}\\[0.3cm]
            \small\textcolor{textPrimary}{Publiczna krytyka, wytykanie błędów, porównywanie i sarkazm.}
        \end{column}
        
        \begin{column}{0.09\textwidth}
            \vspace{1.2cm}
            \centering\textcolor{textMuted}{\Large $\rightarrow$}
        \end{column}
        
        \begin{column}{0.24\textwidth}
            {\fontsize{40}{40}\selectfont\sffamily\bfseries\textcolor{primaryNavy!10}{3}}\\[-0.7cm]
            {\large\bfseries\textcolor{primaryNavy}{Efekt}}\\[0.3cm]
            \small\textcolor{textPrimary}{Opór ucznia, blokada poznawcza, niechęć i rozpad relacji.}
        \end{column}
        
        \begin{column}{0.05\textwidth}
        \end{column}
    \end{columns}
    
    \vspace{1.2cm}
    \pause
    
    \begin{center}
        \normalsize\textcolor{textPrimary}{Używanie wstydu jako narzędzia dydaktycznego}\\
        \vspace{0.15cm}
        \large\textbf{\textcolor{accentCoral}{czynnie niweczy zamierzone cele edukacyjne.}}
    \end{center}
\end{frame}

% --- SLAJD 8: QUIZ 1 ---
\begin{frame}{FAKT CZY MIT?}
    \vfill
    \begin{center}
        {\Large\sffamily\bfseries\textcolor{primaryNavy}{
        ZAWSTYDZENIE WYWOŁUJE POTRZEBĘ\\
        UKRYCIA SIĘ, A NIE CHĘĆ DO ROZWOJU.
        }}\\[1.2cm]
        
        \pause
        {\fontsize{50}{60}\selectfont\sffamily\bfseries\textcolor{accentTeal}{FAKT}}
    \end{center}
    \vfill
\end{frame}

% --- SLAJD 9: QUIZ 2 ---
\begin{frame}{FAKT CZY MIT?}
    \vfill
    \begin{center}
        {\Large\sffamily\bfseries\textcolor{primaryNavy}{
        UCZNIÓW TRZEBA ZASKOCZYĆ PUBLICZNĄ\\
        JEDYNKĄ, WTEDY WEZMĄ SIĘ DO PRACY.
        }}\\[1.2cm]
        
        \pause
        {\fontsize{50}{60}\selectfont\sffamily\bfseries\textcolor{accentCoral}{MIT}}
    \end{center}
    \vfill
\end{frame}

% --- SLAJD 10: QUIZ 3 ---
\begin{frame}{FAKT CZY MIT?}
    \vfill
    \begin{center}
        {\Large\sffamily\bfseries\textcolor{primaryNavy}{
        ODPORNOŚĆ PSYCHICZNĄ BUDUJE SIĘ\\
        W POCZUCIU BEZPIECZEŃSTWA, A NIE STRESU.
        }}\\[1.2cm]
        
        \pause
        {\fontsize{50}{60}\selectfont\sffamily\bfseries\textcolor{accentTeal}{FAKT}}
    \end{center}
    \vfill
\end{frame}

% --- SLAJD 11: QUIZ 4 ---
\begin{frame}{FAKT CZY MIT?}
    \vfill
    \begin{center}
        {\Large\sffamily\bfseries\textcolor{primaryNavy}{
        ŻYCIE NIE JEST KOLOROWE —\\
        ZAWSTYDZAJĄC UCZNIÓW HARTUJĘ ICH NA DOROSŁOŚĆ.
        }}\\[1.2cm]
        
        \pause
        {\fontsize{50}{60}\selectfont\sffamily\bfseries\textcolor{accentCoral}{MIT}}
    \end{center}
    \vfill
\end{frame}

% --- SLAJD 12: QUIZ 5 ---
\begin{frame}{FAKT CZY MIT?}
    \vfill
    \begin{center}
        {\Large\sffamily\bfseries\textcolor{primaryNavy}{
        PORÓWNYWANIE RODZI NIECHĘĆ DO\\
        WZOROWEGO UCZNIA, A NIE MOTYWACJĘ.
        }}\\[1.2cm]
        
        \pause
        {\fontsize{50}{60}\selectfont\sffamily\bfseries\textcolor{accentTeal}{FAKT}}
    \end{center}
    \vfill
\end{frame}

% --- SLAJD 13: DŁUGOFALOWE KONSEKWENCJE ---
{
\usebackgroundtemplate{
    \tikz[remember picture, overlay]
    \node[opacity=0.50, inner sep=0pt] at (current page.center)
    {\includegraphics[width=\paperwidth, height=\paperheight]{ciezar.jpg}};
}
\begin{frame}\normalsize{}
    
    \vspace{0.5cm}
    {\Large\textcolor{textPrimary}{
        Konsekwencje publicznego oceniania, zawstydzania, \\wytykania błędów na forum klasy oraz \\porównywania uczniów między sobą.\\[0.01cm]
        \textbf{\textcolor{accentCoral}{}}
    }}
    
    \vspace{0.04cm}
    \begin{columns}[T]
        \begin{column}{0.10\textwidth}
        \end{column}
        
        \begin{column}{0.50\textwidth}
            \pause
            {\large\sffamily\bfseries\textcolor{primaryNavy}{Wymiar Psychologiczny}}\\[0.2cm]
            \begin{itemize}
                \setlength{\itemsep}{6pt}
                \item \textbf{Wewnętrzny krytyk} 
                \item \textbf{Wyuczona bezradność} 
                \item \textbf{Lęk przed błędem} 
            \end{itemize}
            
            \vspace{0.3cm}
            \pause
            {\large\sffamily\bfseries\textcolor{primaryNavy}{Wymiar Społeczny}}\\[0.2cm]
            \begin{itemize}
                \setlength{\itemsep}{6pt}
                \item \textbf{Rywalizacja zamiast współpracy} 
                \item \textbf{Przemoc, dręczenie} 
                \item \textbf{Niechęć do przedmiotu} 
            \end{itemize}
            
        \end{column}
        
        \begin{column}{0.40\textwidth}
        \end{column}
    \end{columns}
\end{frame}
}

% --- SLAJD 14: POZIOMY KOMUNIKACJI ---
\begin{frame}{Poziomy komunikacji w klasie}
    \vspace{0.4cm}
    \begin{columns}[T]
        \begin{column}{0.04\textwidth}
        \end{column}
        
        \begin{column}{0.28\textwidth}
            {\fontsize{55}{55}\selectfont\sffamily\bfseries\textcolor{accentCoral!15}{1}}\\[-0.8cm]
            {\large\bfseries\textcolor{accentCoral}{Obronna}}\\[0.3cm]
            \footnotesize\textcolor{textPrimary}{\textbf{Niskie zaufanie, strach.}}\\
            Uczeń filtruje przekaz, unika kary i szuka strategii obronnych.\\[0.4cm]
            \normalsize\textbf{\textcolor{textPrimary}{Wynik:}} \textcolor{accentCoral}{Przetrwanie}
        \end{column}
        
        \begin{column}{0.04\textwidth}
            \vspace{1.5cm}
            \centering\textcolor{cardBorder}{\large $\rightarrow$}
        \end{column}
        
        \begin{column}{0.28\textwidth}
            {\fontsize{55}{55}\selectfont\sffamily\bfseries\textcolor{primaryNavy!10}{2}}\\[-0.8cm]
            {\large\bfseries\textcolor{primaryNavy}{Kompromis}}\\[0.3cm]
            \footnotesize\textcolor{textPrimary}{\textbf{Umiarkowane zaufanie.}}\\
            Poszanowanie reguł, ale bez pełnej otwartości na próby i błędy.\\[0.4cm]
            \normalsize\textbf{\textcolor{textPrimary}{Wynik:}} \textcolor{primaryNavy}{Poprawność}
        \end{column}
        
        \begin{column}{0.04\textwidth}
            \vspace{1.5cm}
            \centering\textcolor{cardBorder}{\large $\rightarrow$}
        \end{column}

        \begin{column}{0.28\textwidth}
            {\fontsize{55}{55}\selectfont\sffamily\bfseries\textcolor{accentTeal!15}{3}}\\[-0.8cm]
            {\large\bfseries\textcolor{accentTeal}{Synergiczna}}\\[0.3cm]
            \footnotesize\textcolor{textPrimary}{\textbf{Wysokie zaufanie i relacja.}}\\
            Bezpieczeństwo wyzwala otwartość, dociekliwość i pasję.\\[0.4cm]
            \normalsize\textbf{\textcolor{textPrimary}{Wynik:}} \textcolor{accentTeal}{\textbf{Pełny rozwój}}
        \end{column}
        
        \begin{column}{0.04\textwidth}
        \end{column}
    \end{columns}
    
    \vspace{1.2cm}
    \begin{center}
        \scriptsize\textcolor{textMuted}{Źródło: Model komunikacji Coveya-Stróżyńskiego (B. Niemierko, „Ocenianie szkolne bez tajemnic”, s. 245)}
    \end{center}
\end{frame}

% --- SLAJD 15: PUŁAPKA KOMUNIKATÓW "TY" ---
\begin{frame}{Pułapka komunikatów „TY”}
    \vspace{0.05cm}
    \begin{columns}[T]
        \begin{column}{0.04\textwidth}\end{column}
        \pause
        
        \begin{column}{0.44\textwidth}
            \begin{center}
                {\Large\rmfamily\itshape\bfseries\textcolor{accentCoral}{Przerzucenie winy}}\\[0.1cm]
                
            \end{center}
            \begin{center}
                \footnotesize
                \textcolor{primaryNavy}{\textbf{Frustracja nauczyciela}} \\
                \vspace{0.1cm}
                \textcolor{textMuted}{$\downarrow$} \\
                \vspace{0.1cm}
                \textcolor{accentCoral}{\textit{„Zachowujesz się niegrzecznie!”}} \\
                \vspace{0.1cm}
                \textcolor{textMuted}{$\downarrow$} \\
                \vspace{0.1cm}
                \textcolor{textPrimary}{\textbf{Odbiór:} „Nauczyciel uważa, że jestem zły”}
            \end{center}
        \end{column}
        
        \begin{column}{0.04\textwidth}\end{column}
        \pause
        
        \begin{column}{0.44\textwidth}
            \begin{center}
                {\Large\rmfamily\itshape\bfseries\textcolor{primaryNavy}{Formy zawstydzania}}\\
            \end{center}
            \vspace{0.2cm}
            
            \small\textcolor{textPrimary}{
            Większość takich komunikatów zawiera oskarżycielskie \textbf{„TY”} i uderza bezpośrednio w poczucie wartości ucznia:
            }
            
            \vspace{0.15cm}
            \begin{itemize}
                \footnotesize
                \setlength{\itemsep}{4pt}
                \item \textbf{Krytykowanie:} \textit{„Nie myślisz dojrzale.”}
                \item \textbf{Epitety:} \textit{„Zachowujesz się jak dziecko!”}
                \item \textbf{Sarkazm:} \textit{„Drugi Einstein z ciebie...”}
                \item \textbf{Moralizowanie:} \textit{„Powinieneś to umieć!”}
            \end{itemize}
        \end{column}
        
        \begin{column}{0.04\textwidth}\end{column}
    \end{columns}
    
    \vspace{0.5cm}
    \pause
    \begin{center}
        \footnotesize \textbf{\textcolor{accentCoral}{Wniosek:}} Komunikat „TY” unika odpowiedzialności za uczucia nauczyciela.\\Uczeń jest wystawiony na publiczną, negatywną ocenę.
        
        \vspace{0.25cm}
        \scriptsize\textcolor{textMuted}{Źródło: Thomas Gordon, \textit{Wychowanie bez porażek w szkole}, s. 150–152}
    \end{center}
\end{frame}

% --- SLAJD 16: ZMIANA JĘZYKA ---
\begin{frame}{Dobór słów ma znaczenie}
    \vspace{0.6cm}
    \begin{columns}[T]
        \begin{column}{0.04\textwidth}
        \end{column}
        
        \begin{column}{0.44\textwidth}
            \begin{center}
                {\Large\rmfamily\itshape\bfseries\textcolor{accentCoral}{Komunikat „TY”}}\\[0.1cm]
                {\footnotesize\textcolor{textMuted}{(oceniający)}}
            \end{center}
            \vspace{0.4cm}
            
            \begin{center}
                \textit{\small „Jesteś nieuważny i znowu przeszkadzasz w prowadzeniu zajęć.”}
            \end{center}
            \vspace{0.4cm}
            
            \footnotesize \textbf{\textcolor{accentCoral}{Skutek:}} Koncentracja na wadach ucznia. Wywołuje natychmiastowy opór i defensywę.
        \end{column}
        
        \begin{column}{0.04\textwidth}
        \end{column}
        
        \begin{column}{0.44\textwidth}
            \begin{center}
                {\Large\rmfamily\itshape\bfseries\textcolor{accentTeal}{Komunikat „JA”}}\\[0.1cm]
                {\footnotesize\textcolor{textMuted}{(opisowy)}}
            \end{center}
            \vspace{0.4cm}
            
            \begin{center}
                \textit{\small „Trudno mi kontynuować wywód, gdy w klasie panuje hałas. Zależy mi na sprawnym przejściu przez ten materiał.”}
            \end{center}
            \vspace{0.15cm}
            
            \footnotesize \textbf{\textcolor{accentTeal}{Skutek:}} Skupienie na faktach i potrzebach. Nie podważa godności odbiorcy.
        \end{column}
        
        \begin{column}{0.04\textwidth}
        \end{column}
    \end{columns}
\end{frame}

% --- SLAJD 17: ZASADA ZERO (POUFNOŚĆ) ---
\begin{frame}{Poufność oceniania}
    \vspace{0.4cm}
    \begin{columns}[T]
        \begin{column}{0.55\textwidth}
        \pause
            \begin{tcolorbox}[moderncard]
                \small
                
                \begin{itemize}
                    \item \textbf{Według RODO} oceny to chronione dane osobowe.
                    \item Oceny są jawne wyłącznie dla osób uprawnionych (uczeń, rodzice, nauczyciel).
                    \item Nie jest wskazane wyczytywanie ocen na głos czy publikacja list wyników, gdy możliwa jest identyfikacja ucznia.
                \end{itemize}
            \end{tcolorbox}
        \end{column}
        \pause
        \begin{column}{0.40\textwidth}
            \begin{tcolorbox}[tealcard]
                \centering
                \textbf{\textcolor{accentTeal}{Dobra praktyka}}\\[0.3cm]
                \small
                Informację zwrotną przekazuj \textbf{dyskretnie i indywidualnie}.
            
            \end{tcolorbox}
        \end{column}
    \end{columns}
\end{frame}

% --- SLAJD 18: OCENIANIE KSZTAŁTUJĄCE ---
{
\usebackgroundtemplate{
    \tikz[remember picture, overlay]
    \node[anchor=east, opacity=0.50, inner sep=0pt] at (current page.east) 
    {\includegraphics[height=\paperheight, keepaspectratio]{schody.jpg}};
}
\begin{frame}{Ocenianie kształtujące w praktyce}
    \vspace{0.4cm}
    \begin{columns}[T]
        \begin{column}{0.55\textwidth}
            \pause
            {\large\sffamily\bfseries\textcolor{primaryNavy}{Informacja zwrotna:}}\\[0.2cm]
            \small
            \begin{enumerate}
                \setlength{\itemsep}{5pt}
                \item Wskazanie elementów wykonanych poprawnie.
                \item Określenie miejsca powstania błędu.
                \item Sugestia konkretnych działań naprawczych.
                \item Wyznaczenie kolejnego celu rozwojowego.
            \end{enumerate}
            
            \vspace{0.5cm}
            \pause
            \scriptsize\textcolor{textMuted}{
                Podawanie uczniowi samego komentarza orientuje go \\na zadanie i skutkuje przyrostem osiągnięć, podczas \\gdy podawanie oceny i komentarza sprawia, że uczeń \\skupia się wyłącznie na cyfrze. \\- B. Niemierko, \textit{„Ocenianie szkolne bez tajemnic”}, s. 253
            }
            
        \end{column}
        
        \begin{column}{0.45\textwidth}
        \end{column}
    \end{columns}
\end{frame}
}

% --- SLAJD 19: PODSUMOWANIE ---
{
\usebackgroundtemplate{
    \tikz[remember picture, overlay]
    \node[anchor=east, opacity=0.50, inner sep=0pt] at (current page.east) 
    {\includegraphics[height=\paperheight, keepaspectratio]{schody.jpg}};
}
\begin{frame}{Podsumowanie}
    \vspace{0.4cm}
    \begin{columns}[T]
        \begin{column}{0.58\textwidth}
            
            \pause
            \small\textcolor{textPrimary}{
            Rola nauczyciela wykracza poza \\przekazywanie wiedzy:
            }
            
            \vspace{0.3cm}
            
            \begin{itemize}
                \setlength{\itemsep}{10pt} 
                \pause\item Sposób komunikacji kształtuje postawę \\ucznia wobec nauki na całe życie.
                \pause\item Publiczne ocenianie i zawstydzanie \\budują bariery rozwojowe.
                \pause\item Poczucie bezpieczeństwa i szacunek \\stanowią fundament efektywnej edukacji.
            \end{itemize}
            
        \end{column}
        
        \begin{column}{0.38\textwidth}
        \end{column}
    \end{columns}
\end{frame}
}

% --- SLAJD 20: ZAKOŃCZENIE ---
\begin{frame}
    \vfill

    \begin{tcolorbox}[moderncard]
        \centering

        \vspace{0.5cm}

        {\Huge\bfseries
        \textcolor{primaryNavy}{Dziękujemy za uwagę!}
        \par}
        \vspace{0.5cm}

        {\large
        \textcolor{primaryNavy}{Jolanta Sułkowska}
        \par}

        \vspace{0.2cm}

        {\large
        \textcolor{primaryNavy}{Karolina Studzińska}
        \par}

        \vspace{0.5cm}
    \end{tcolorbox}

    \vfill
\end{frame}
\end{document}
